\chapter{Conclusiones y Líneas Futuras}
\label{cap:conclusiones}

Este Trabajo de Fin de Grado se propuso abordar los desafíos inherentes a la imagen cuantitativa con el isótopo no convencional Itrio-86 ($^{86}\text{Y}$), utilizando un enfoque sistemático basado en simulación Monte Carlo. A través de la validación del entorno de simulación, la caracterización de los artefactos y la optimización de parámetros, se ha desarrollado una metodología robusta que culmina en una exitosa prueba de concepto para la imagen PET multiplexada.

\section{Conclusiones}
\label{sec:conclusiones:conclusiones}

El primer hito fundamental de este trabajo ha sido la validación exitosa de un entorno de simulación completo para isótopos complejos. Se ha demostrado que la combinación del framework \texttt{penRed}, junto con su módulo \texttt{PENNUC} y un software de post-procesado desarrollado a medida, es capaz de reproducir fielmente la compleja física de desintegración del $^{86}\text{Y}$. Esta capacidad incluye la simulación precisa de la emisión de \textit{prompt gammas} y la generación de coincidencias triples, lo cual ha sido corroborado por la excelente concordancia entre el espectro energético simulado y los datos nucleares de referencia (ENSDF).

Una vez validada la herramienta, el estudio permitió cuantificar la degradación de la imagen y aislar el impacto específico de los \textit{prompt gammas}. La metodología incremental reveló que la combinación de efectos físicos, como la atenuación y el \textit{scatter}, sumada a una respuesta electrónica realista con coincidencias aleatorias, provoca una pérdida de contraste (CRC) superior al \SI{93}{\percent} en comparación con un escenario ideal. La comparativa directa con un isótopo \enquote{limpio} como el $^{18}\text{F}$ probó de manera concluyente que la fuente primaria de esta degradación no es el \textit{scatter} convencional, sino el masivo fondo de ruido inducido por las cascadas gamma del $^{86}\text{Y}$.

En respuesta a esta problemática, la optimización paramétrica identificó que el ajuste de la Ventana de Energía (EW) es la estrategia más eficaz para la recuperación del contraste. El barrido experimental demostró que una configuración con una EW del \SI{10}{\percent} es óptima, logrando no solo recuperar la señal perdida, sino mejorar el contraste en un \SI{63}{\percent} respecto a la línea base realista, manteniendo al mismo tiempo una excelente Relación Señal-Ruido (SNR). Por otro lado, se observó que la Ventana de Coincidencia (CW) tiene un impacto secundario, contribuyendo principalmente a la mejora de la SNR mediante la reducción de aleatorias.

Finalmente, se ha demostrado la viabilidad técnica de la imagen PET multiplexada. La prueba de concepto ha sido exitosa gracias al algoritmo de post-procesado implementado, que logró identificar y aislar las coincidencias triples para generar un \textit{dataset} exclusivo del $^{86}\text{Y}$. La posterior sustracción de esta señal permitió recuperar la distribución espacial del $^{18}\text{F}$ con una contaminación cruzada mínima, validando así el principio de utilizar los \textit{prompt gammas} no como una fuente de ruido, sino como una firma física única para la separación de isótopos en adquisiciones simultáneas.

\section{Líneas Futuras}
\label{sec:conclusiones:lineas_futuras}


